% -*-latex-*-
%Farmers and thieves
\begin{figure}
\begin{center}
\mbox{%
\beginpicture
\setcoordinatesystem units <3in,10in>
\setplotarea x from 0 to 1, y from 1 to 1.2
\axis left 
  label {Fitness}
  ticks withvalues {IL} {SL} / at 1.1 1.2 / /
\axis bottom 
  label {Frequency of SL}
  ticks numbered from 0 to 1 by 1 /
\axis top ticks withvalues {ESS} / at 0.556 /  /
\axis right /
% Cultural Evolution model from Rogers (1988)
%
% > wS := w + b*(1-p)*(1-u)/( 1 - p*(1-u));
%
%                                     b (1 - p) (1 - u)
%                           wS := w + -----------------
%                                       1 - p (1 - u)
% 
% > wi := w + b*(1-c);
%                                    wi := 1.1
% > wI := w + b*(1-c);
%                               wI := w + b (1 - c)
%
% Parameter values: w = b = 1, c = 0.9, s = 0, and u = 0.8.

\put {1} [t] <10pt,-6pt> at 0 1.1
\put {2} [t] <0pt,-10pt> at 0.05 1.191919192
\put {3} [b] <0pt,10pt> at 0.95 1.012345679

\plot 
0 1.2 0.05 1.191919192 0.10 1.183673469 0.15 1.175257732
0.20 1.166666667 0.25 1.157894737 0.30 1.148936170 0.35 1.139784946
0.40 1.130434783 0.45 1.120879121 0.50 1.111111111 0.55 1.101123596
0.60 1.090909091 0.65 1.080459770 0.70 1.069767442 0.75 1.058823529
0.80 1.047619048 0.85 1.036144578 0.90 1.024390244 0.95 1.012345679
1.00 1
/
% Fitness of Individual Learner
\plot 0 1.1  1 1.1 /
\setdots
\putrule from 0.556 1 to 0.556 1.2
\setsolid
\put {$\Leftarrow$} [l] at 0.556 1.15
\put {$\Rightarrow$} [r] at 0.556 1.15
\endpicture%
}
\end{center}
\caption{The fitnesses of both learning strategies and also the mean
  fitness of the population are plotted against the frequency of social
  learning, assuming that $w = b = 1$, $c = 0.9$, $s = 0$, and $u = 0.8$.  The
  numbers refer to the following assumptions: (1) the costs and benefits of
  individual learning are independent of its frequency; (2) when social
  learning is rare, its fitness value exceeds that of individual learning
  because it provides relatively up-to-date information at low cost; and (3)
  when social learning is very common, its fitness value falls below that of
  individual learning because the information it provides is out of date.
\label{fig.cult}}
\end{figure}

%%% Local Variables: 
%%% mode: latex
%%% TeX-master: "sl"
%%% End: 
